% Generated draft: requires mathematical and source review.
\documentclass[11pt]{article}
\usepackage[margin=25mm]{geometry}
\usepackage{fontspec}
\setmainfont{Latin Modern Roman}
\usepackage{amsmath,amssymb,mathtools,unicode-math}
\setmathfont{Latin Modern Math}
\usepackage{enumitem,needspace,xurl,hyperref}
\setlength{\parindent}{0pt}
\setlength{\parskip}{5pt}
\setlength{\emergencystretch}{3em}
\newcommand{\N}{\mathbb{N}}
\newcommand{\Z}{\mathbb{Z}}
\newcommand{\Q}{\mathbb{Q}}
\newcommand{\R}{\mathbb{R}}
\newcommand{\C}{\mathbb{C}}
\newcommand{\F}{\mathbb{F}}
\newcommand{\A}{\mathbb{A}}
\newcommand{\PP}{\mathbb{P}}
\newcommand{\E}{\mathbb{E}}
\newcommand{\eps}{\varepsilon}
\newcommand{\dd}{\,\mathrm{d}}
\begin{document}
\section*{UnsolvedMath ID: 30005129}
\textbf{Source problem code:} catalogue entry
\subsection{Definitions and mathematical statement}
For a small category \(\mathcal C\), its classifying space is \(B\mathcal C=|N_\bullet\mathcal C|\), the geometric realization of its nerve. Homology of a category can thus mean homology of its classifying space. The supplied record does not define the reductive Borel--Serre categories under consideration or specify their underlying algebraic or geometric data.

To formulate homological stability, one must specify a family of categories \((\mathcal C_n)\), stabilization functors \(s_n:\mathcal C_n\to\mathcal C_{n+1}\), and coefficients. For example, with a fixed abelian group \(A\) as constant coefficients, a stability statement would specify degrees \(i\) and indices \(n\) for which
\[
(Bs_n)_*:H_i(B\mathcal C_n;A)\longrightarrow H_i(B\mathcal C_{n+1};A)
\]
is an isomorphism, together with any separate surjectivity range. This describes the form of a possible statement, not a stability assertion supplied by the record; neither the family, stabilization functors, coefficients, nor intended ranges are given.

\textbf{Source question.} Determine the homological-stability properties of the intended reductive Borel--Serre categories, and determine whether the intended associated structure possesses higher homotopy commutativity. The original passage additionally asks, more vaguely, about exploiting their geometric origins.

Higher homotopy commutativity requires a specified multiplication and a specified notion of coherent commutativity. Neither is supplied, and the antecedent of ``it'' in the original question is unclear. In particular, the record does not identify whether this property concerns individual classifying spaces, a construction assembling the categories, or another associated object. Consequently, the supplied evidence determines a research question but does not determine a unique fully quantified conjecture. No particular stability range or coherence structure is asserted here.
\subsection{Short English statement}
The problem asks how homology stabilizes for reductive Borel--Serre categories and whether an associated multiplication is commutative with suitable higher homotopies. It also asks how their geometric origins might be exploited. The categories, stabilization maps, coefficients, multiplication, and required coherence are unspecified, so a precise conjecture requires further source clarification.
\subsection{Sources}
\begin{itemize}
\item Supplied local UnsolvedMath registry record, permanent ID \(30005129\), titled ``Homological Stability of Reductive Borel--Serre Categories''; exact source problem code \texttt{catalogue entry}. No individual proposer is identified; the supplied proposed year is \(2022\).
\item Supplied report citation: \emph{Algebraic K-Theory (2023)}, Oberwolfach Reports, DOI \texttt{10.4171/owr/2022/24}, pp.\ 1313--1376. \url{https://doi.org/10.4171/owr/2022/24}. The record attributes the problem to workshop open-problem discussions and retains an extracted report passage as context.
\item Randal-Williams, ``RBS stability'' (supplied reference label): \url{https://arxiv.org/abs/2405.07566}.
\item Jansen, ``unstable K-theory'' (supplied reference label): \url{https://arxiv.org/abs/2405.02065}.
\item Clausen--Jansen, ``RBS model'' (supplied reference label): \url{https://doi.org/10.1007/s00029-023-00900-8}.
\item Literature assessment dated exactly \texttt{2026-08-22}: the supplied curation labels the problem partially solved and reports substantial proved stability ranges for reductive Borel--Serre spaces and related unstable algebraic K-theory, while describing the broader requested theory as incomplete. This is supplied dated curation, not an independently verified assessment; the record requests independent verification of time-sensitive status and cited claims.
\item The registry labels the statement \texttt{corrected\_verified} and reports that the official report was checked directly. These are supplied provenance claims; no independent source inspection is asserted in this draft.
\end{itemize}
\end{document}
